% Chapter 6: the planned structure of the final thesis.
\chapter{Thesis chapter plan}\label{ch:plan}

The thesis follows a monograph structure with three core chapters, one
per research question, each written to stand as a journal article.
Table~\ref{tab:chapters} lists the chapters with their status at the
time of this review.

\begin{table}[ht]
\centering\small
\caption{Planned chapters, their questions and their status.}
\label{tab:chapters}
\begin{tabularx}{\linewidth}{@{}>{\sffamily\color{accent}}l X
  >{\sffamily\footnotesize}l@{}}
\toprule
\thead{No.} & \thead{Chapter and content} & \thead{Status}\\
\midrule
1 & \textbf{Introduction.} Heat exposure, sensing gaps, aims and
    thesis outline. & outline \\
2 & \textbf{Background.} Urban climate mapping, Gaussian processes,
    adaptive design. & draft \\
3 & \textbf{Data and model.} Aldmere network, covariates, the model of
    Eq.~\eqref{eq:model} and its validation. & draft \\
4 & \textbf{Exposure-weighted placement.} The rule of
    Eq.~\eqref{eq:acq}, theory and simulations; answers \ref{rq:acq}.
    & in progress \\
5 & \textbf{Learning across cities.} Transfer to Lyre and Westhaven;
    answers \ref{rq:transfer}. & planned \\
6 & \textbf{A season on the street.} Field campaign and its lessons;
    answers \ref{rq:field}. & planned \\
7 & \textbf{Conclusions.} Synthesis, guidance for city climate
    offices, limitations and future work. & planned \\
\bottomrule
\end{tabularx}
\end{table}

Chapter~4 is planned as a submission to \emph{Environmental Modelling
\& Software} in spring 2027, Chapter~5 to \emph{Urban Climate}, and
Chapter~6 to \emph{Nature Cities} or a comparable venue, depending on
the strength of the field result.
